\documentclass[10pt]{article}

\usepackage[a4paper,margin=0.9in]{geometry}
\usepackage{amsmath,amssymb,amsthm,mathtools}
\usepackage{microtype}
\usepackage{xcolor}
\usepackage[colorlinks=true,linkcolor=blue!50!black,
  citecolor=blue!50!black,urlcolor=blue!50!black]{hyperref}

\newtheorem{theorem}{Theorem}[section]
\newtheorem{proposition}[theorem]{Proposition}
\newtheorem{lemma}[theorem]{Lemma}
\newtheorem{corollary}[theorem]{Corollary}
\theoremstyle{definition}
\newtheorem{definition}[theorem]{Definition}

\newcommand{\C}{\mathbb C}
\newcommand{\R}{\mathbb R}
\newcommand{\Z}{\mathbb Z}
\newcommand{\N}{\mathbb N}
\newcommand{\T}{\mathbb T}
\newcommand{\HS}{\mathrm{HS}}
\newcommand{\proj}{\mathrm{proj}}
\newcommand{\Hwater}{\mathcal H}
\newcommand{\Bhaar}{\mathcal B_{\mathrm H}}
\newcommand{\ee}{\mathrm e}

\title{\textbf{Supplementary Information for\\
``A universal information law for quantum return''}}
\author{Michael Haufschild}
\date{28 July 2026}

\begin{document}
\maketitle

\section{Definitions and exact projective geometry}

Let \(w=(w_1,\ldots,w_r)\) be positive with \(\sum_jw_j=1\), and define
\[
 A_w(n;\omega)=\left|\sum_jw_j\ee^{in\omega_j}\right|.
\]
Adding the same scalar to all \(\omega_j\) leaves \(A_w\) unchanged.
Accordingly, one coordinate may be used as a projective pivot.

\begin{proposition}[Weighted projective variance]
\label{prop:variance}
For unit phases \(z_j\in\T\),
\begin{equation}
 \sum_{j,k}w_jw_k|z_j-z_k|^2
 =2\left(1-\left|\sum_jw_jz_j\right|^2\right).
\label{eq:variance}
\end{equation}
For \(\theta\in\R^r\),
\begin{equation}
 \frac12\sum_{j,k}w_jw_k(\theta_j-\theta_k)^2
 =\sum_jw_j\theta_j^2-\left(\sum_jw_j\theta_j\right)^2.
\label{eq:real-variance}
\end{equation}
\end{proposition}

\begin{proof}
Expand the squared differences, use \(|z_j|=1\), and factor the cross
term as
\(\sum_{j,k}w_jw_kz_j\overline z_k=|\sum_jw_jz_j|^2\).
The same expansion over real scalars proves \eqref{eq:real-variance}.
\end{proof}

The covariance term on the right of \eqref{eq:real-variance} shows that the
separable water program used in the theorem is not being identified with the
exact quotient metric.  It is the common anisotropic envelope arising in the
covering and target-volume arguments.

\section{Reverse water filling}

For a pivot \(k\), define
\[
 \Hwater_k(w,\eta)
 =\min\left\{\frac12\sum_{j\ne k}\log\frac{w_j}{D_j}:
   0<D_j\leq w_j,\quad\sum_{j\ne k}D_j\leq\eta\right\}.
\]

\begin{proposition}[Global optimizer]
\label{prop:optimizer}
If \(\eta\geq1-w_k\), then \(\Hwater_k(w,\eta)=0\).  If
\(0<\eta<1-w_k\), there is an active water level
\(\lambda_k>0\) satisfying
\[
 \sum_{j\ne k}\min\{w_j,\lambda_k^{-1}\}=\eta,
\]
and
\[
 D_j^*=\min\{w_j,\lambda_k^{-1}\},\qquad
 \Hwater_k(w,\eta)=\frac12\sum_{j\ne k}\log_+(\lambda_kw_j).
\]
\end{proposition}

\begin{proof}
Set \(x_j=w_j/D_j\).  The feasible conditions are
\[
 x_j\geq1,\qquad \sum_{j\ne k}\frac{w_j}{x_j}\leq\eta,
\]
and the objective is \(\frac12\sum_{j\ne k}\log x_j\).
When \(\eta\geq1-w_k\), \(x_j=1\) is feasible and has the smallest possible
objective.

Suppose \(\eta<1-w_k\).  The continuous nonincreasing function
\[
 f(\lambda)=\sum_{j\ne k}\min\{w_j,\lambda^{-1}\}
\]
starts above \(\eta\) and tends to zero.  On the relevant range at least one
coordinate is active.

For any \(a\geq0\), \(\lambda>0\), and \(x\geq1\),
\begin{equation}
 \log\max\{1,\lambda a\}
 +\lambda\left(\min\{a,\lambda^{-1}\}-\frac{a}{x}\right)
 \leq\log x.
\label{eq:support}
\end{equation}
If \(\lambda a\leq1\), this is \(1-x^{-1}\leq\log x\) after
multiplication by \(\lambda a\leq1\).  If \(\lambda a\geq1\), it is
\(\log y\leq y-1\) with \(y=x/(\lambda a)\), after rearrangement.
Sum \eqref{eq:support} over the non-pivot coordinates.  The water equation
and feasibility make the correction term nonnegative, proving global
optimality.  Evaluating the objective gives the stated logarithmic sum.
\end{proof}

\section{Universal and Haar-quantile recurrence laws}

\subsection{Constructive ceiling}

\begin{lemma}[Anisotropic product-box return]
\label{lem:box}
Let \(x_j\geq1\), \(j\ne k\), satisfy
\(\sum_{j\ne k}w_j/x_j\leq\eta\).  For every
\(\omega\in\R^r\), there is \(1\leq n\leq P\), where
\[
 P=\prod_{j\ne k}\left\lceil\pi\sqrt{2x_j}\right\rceil,
\]
such that \(A_w(n;\omega)\geq1-\eta\).
\end{lemma}

\begin{proof}
Put \(Q_j=\lceil\pi\sqrt{2x_j}\rceil\).  Partition the \(j\)th circle into
\(Q_j\) half-open arcs.  Among the \(P+1\) points
\[
 \bigl(m(\omega_j-\omega_k)\bmod2\pi\bigr)_{j\ne k},
 \qquad 0\leq m\leq P,
\]
two occupy the same product box.  Subtract their indices.  This gives
\(1\leq n\leq P\) and integers \(q_j\) with
\[
 |n(\omega_j-\omega_k)-2\pi q_j|\leq\frac{2\pi}{Q_j}.
\]
Set the pivot error to zero and call the remaining centered errors
\(\theta_j\).  Then
\begin{align*}
 A_w(n;\omega)
 &\geq\Re\sum_jw_j\ee^{i\theta_j}\\
 &\geq1-\frac12\sum_{j\ne k}w_j\theta_j^2\\
 &\geq1-2\pi^2\sum_{j\ne k}\frac{w_j}{Q_j^2}\\
 &\geq1-\sum_{j\ne k}\frac{w_j}{x_j}
 \geq1-\eta.
\end{align*}
\end{proof}

At the water optimizer,
\(\lceil\pi\sqrt{2x_j}\rceil\leq2\pi\sqrt{2x_j}\).  Thus
\[
 \log T_w(\eta)
 \leq\Hwater_{\proj}(w,\eta)+(r-1)\log(2\pi\sqrt2).
\]

\begin{lemma}[Selected-support product-box return]
\label{lem:selected-box}
Fix a pivot \(k\), a set
\(S\subseteq\{1,\ldots,r\}\setminus\{k\}\), and \(x_j\geq1\).
Put
\[
 \tau_S=\sum_{j\ne k,\,j\notin S}w_j.
\]
If
\[
 2\tau_S+\sum_{j\in S}\frac{w_j}{x_j}\leq\eta,
\]
then every phase packet has a return \(1\leq n\leq P_S\), where
\[
 P_S=\prod_{j\in S}\left\lceil\pi\sqrt{2x_j}\right\rceil,
\]
and
\[
 \log P_S\leq\frac12\sum_{j\in S}\log x_j
 +|S|\log(2\pi\sqrt2).
\]
\end{lemma}

\begin{proof}
Give a selected coordinate
\(Q_j=\lceil\pi\sqrt{2x_j}\rceil\) phase cells and every other coordinate
one cell.  The anisotropic product-box argument gives \(1\leq n\leq P_S\)
and errors \(|\theta_j|\leq2\pi/Q_j\) for \(j\in S\).  After removing the
pivot phase, use \(\cos\theta_j\geq1-\theta_j^2/2\) on \(S\) and
\(\cos\theta_j\geq-1\) off \(S\):
\[
\begin{aligned}
 A_w(n;\omega)
 &\geq w_k+\sum_{j\in S}w_j(1-\theta_j^2/2)
        -\sum_{j\ne k,\,j\notin S}w_j\\
 &\geq1-2\tau_S-\sum_{j\in S}\frac{w_j}{x_j}
 \geq1-\eta.
\end{aligned}
\]
Only selected coordinates occur in the product.  Applying
\(\lceil\pi\sqrt{2x}\rceil\leq2\pi\sqrt{2x}\) to them proves the logarithmic
bound.
\end{proof}

\subsection{Two-sided target measure, Haar unions and universal-deadline floor}

Let normalized Haar measure on \(\T^r\) be denoted by \(m_{\T^r}\), and set
\[
 \mu_w(\eta)=m_{\T^r}
 \left\{z:\left|\sum_jw_jz_j\right|\geq1-\eta\right\}.
\]

\begin{lemma}[Constructive inner box]
\label{lem:inner-box}
For every pivot \(k\) and every feasible scale \(x_j\geq1\) with
\(\sum_{j\ne k}w_j/x_j\leq\eta\),
\[
 \mu_w(\eta)\geq
 \prod_{j\ne k}\frac{\sqrt{2/x_j}}{\pi}
 =\left(\frac{\sqrt2}{\pi}\right)^{r-1}
   \exp\{-\Hwater_k(x)\},
\]
where \(\Hwater_k(x)=\frac12\sum_{j\ne k}\log x_j\).
\end{lemma}

\begin{proof}
Apply the Haar-preserving projective shear
\[
 (z_j)_j\longmapsto
 \left(z_k,(z_j/z_k)_{j\ne k}\right).
\]
In the relative coordinates \(y_j\in\mathbb R/\mathbb Z\), take the product box
\[
 \|y_j\|_{\R/\Z}\leq\frac{\sqrt{2/x_j}}{2\pi}.
\]
The circle chord satisfies
\(|\ee^{2\pi iy_j}-1|\leq2\pi\|y_j\|_{\R/\Z}\), hence the weighted pivot
energy is at most \(2\sum_{j\ne k}w_j/x_j\leq2\eta\).  For arbitrary unit
phases,
\[
 1-\frac12\sum_jw_j|z_j-z_k|^2
 \leq\left|\sum_jw_jz_j\right|,
\]
so the box lies in the coherence target.  Its pivot coordinate has measure
one; each other interval has normalized measure
\(\sqrt{2/x_j}/\pi\).  Multiplication gives the lower bound, and
\(\sqrt{2/x}/\pi=(\sqrt2/\pi)\exp\{-\frac12\log x\}\) gives the second.
\end{proof}

\begin{lemma}[Pivot-resolved weighted target volume]
\label{lem:target}
\[
 \mu_w(\eta)\leq\Bhaar(w,\eta):=
 \sum_k\exp\left\{\frac{\pi}{4}\lambda_k\eta
 -\Hwater_k(w,\eta)\right\}.
\]
\end{lemma}

\begin{proof}
On the target event, Proposition \ref{prop:variance} gives
\[
 \sum_{j,k}w_jw_k|z_j-z_k|^2\leq4\eta.
\]
After averaging over \(k\), at least one pivot satisfies
\[
 \sum_jw_j|z_j-z_k|^2\leq4\eta.
\]
The projective shear \((z_k,(z_j/z_k)_{j\ne k})\) preserves Haar measure.
For
\[
 M(u)=\int_{\T}\exp(-u|z-1|^2)\,dz,
\]
exponential Markov gives, for every \(s>0\),
\[
 m_{\T^r}\left\{\sum_jw_j|z_j-z_k|^2\leq4\eta\right\}
 \leq\ee^{4s\eta}\prod_{j\ne k}M(sw_j).
\]
On the standard angular interval,
\(|\ee^{i\theta}-1|\geq2|\theta|/\pi\).  Extending the resulting Gaussian
integral to \(\R\) gives
\[
 M(u)\leq\min\left\{1,\frac{\sqrt\pi}{4\sqrt u}\right\}.
\]
For an active pivot choose \(s=\pi\lambda_k/16\).  The product bound becomes
\[
 \exp\left\{\frac{\pi}{4}\lambda_k\eta-\Hwater_k(w,\eta)\right\}.
\]
For an inactive pivot take \(\lambda_k=1\); the same envelope follows from
the coordinate moment bound.  Union over the pivot events gives the displayed
sum without replacing pivot-specific scales by full support size.
\end{proof}

\begin{theorem}[Two-sided target-information law]
\label{thm:target-information}
For \(r\geq2\) and \(0<\eta\leq1\), put
\(I_w(\eta)=-\log\mu_w(\eta)\).  Then
\[
 \Hwater_{\proj}-\frac{\pi}{4}(r-1)-\log r
 \leq I_w
 \leq\Hwater_{\proj}-(r-1)\log\frac{\sqrt2}{\pi}.
\]
In particular \(\mu_w(\eta)>0\).  If \(\Hwater_{\proj}>0\), then
\[
 \left|\frac{I_w}{\Hwater_{\proj}}-1\right|
 \leq
 \frac{\max\{\frac{\pi}{4}(r-1)+\log r,\,
 -(r-1)\log(\sqrt2/\pi)\}}{\Hwater_{\proj}}.
\]
\end{theorem}

\begin{proof}
Use Lemma \ref{lem:inner-box} at a minimizing canonical water scale.  This
gives
\[
 \left(\frac{\sqrt2}{\pi}\right)^{r-1}
 \ee^{-\Hwater_{\proj}}\leq\mu_w(\eta).
\]
The canonical water levels satisfy \(\lambda_k\eta\leq r-1\), while
\(\Hwater_{\proj}\leq\Hwater_k\).  Lemma \ref{lem:target} therefore gives
\[
 \mu_w(\eta)\leq
 r\exp\{\tfrac{\pi}{4}(r-1)-\Hwater_{\proj}\}.
\]
Both bounding expressions are positive.  Apply the increasing logarithm,
reverse signs, and rearrange.  The absolute and relative forms follow by
taking the larger one-sided remainder.
\end{proof}

\begin{theorem}[Finite-horizon Haar quantile]
\label{lem:avoid}
Let
\[
 E_N=\{\omega:\exists\,1\leq n\leq N,\,
 A_w(n;\omega)\geq1-\eta\}.
\]
Then
\[
 m_{\T^r}(E_N)\leq N\Bhaar(w,\eta).
\]
In particular, if \(N\Bhaar(w,\eta)\leq\alpha\), the complement has Haar
measure at least \(1-\alpha\).
\end{theorem}

\begin{proof}
For each positive \(n\), the componentwise power map \(z\mapsto z^n\)
preserves Haar measure.  The recurrence target at time \(n\) therefore has
measure \(\mu_w(\eta)\).  The union of the first \(N\) targets has measure at
most \(N\mu_w(\eta)\leq N\Bhaar(w,\eta)\).  Taking complements proves the
last assertion.
\end{proof}

The active water equation implies
\[
 \lambda_k\eta
 =\sum_{j\ne k}\min\{\lambda_kw_j,1\}\leq r-1.
\]
The inactive convention gives
\(\lambda_k\eta=\eta\leq1\leq r-1\).  Since
\(\Hwater_{\proj}\leq\Hwater_k\),
\[
 \Bhaar(w,\eta)\leq
 r\exp\left\{\frac{\pi}{4}(r-1)-\Hwater_{\proj}(w,\eta)\right\}.
\]

Set
\[
 B=\frac{1}{2r}
 \exp\left\{\Hwater_{\proj}(w,\eta)-\frac{\pi}{4}(r-1)\right\}.
\]
If \(B<1\), then \(T_w(\eta)\geq1>B\).  Otherwise take
\(N=\lfloor B\rfloor\); Lemmas \ref{lem:target} and \ref{lem:avoid} give
\(T_w(\eta)>N\), hence the integer \(T_w(\eta)\geq N+1>B\).
Equivalently,
\[
 \Hwater_{\proj}(w,\eta)-\frac{\pi}{4}(r-1)-\log(2r)
 <\log T_w(\eta).
\]
Together with the constructive ceiling, this proves the main
recurrence--rate-distortion sandwich.

\begin{corollary}[Return after a prescribed delay]
\label{cor:post-delay}
For every compatible phase packet \(\omega\) and every \(M\in\N\), there is
an integer \(n\) such that
\[
 M<n\leq(M+1)T_w(\eta),\qquad A_w(n;\omega)\geq1-\eta.
\]
In particular, if the packet has exited the target by tick \(M\), this is a
post-exit recurrence.
\end{corollary}

\begin{proof}
Apply the definition of \(T_w(\eta)\) to the compatible packet
\((M+1)\omega\).  It supplies \(1\leq q\leq T_w(\eta)\) with
\(A_w(q;(M+1)\omega)\geq1-\eta\).  Put \(n=(M+1)q\).  Then \(M<n\),
\(n\leq(M+1)T_w(\eta)\), and
\(A_w(n;\omega)=A_w(q;(M+1)\omega)\).
\end{proof}

\section{Finite observer-coherence theorem}

Let \(z_i\in\T\), let \(B^{(a)}\in\C^{d\times d}\), and let
\(\lambda_a\geq0\).  Define
\[
 W_{ij}=\sum_a\lambda_a|B^{(a)}_{ij}|^2,\qquad
 D^2(z)=\sum_a\lambda_a
 \|\operatorname{diag}(z)B^{(a)}
   \operatorname{diag}(z)^\dagger-B^{(a)}\|_{\HS}^2.
\]

\begin{theorem}[Exact phase energy and zero set]
\label{thm:observer}
\[
 D^2(z)=\sum_{i,j}W_{ij}|z_i\overline z_j-1|^2.
\]
If \(G\) has an edge \(\{i,j\}\) precisely when
\(W_{ij}>0\) or \(W_{ji}>0\), then
\[
 D^2(z)=0
 \quad\Longleftrightarrow\quad
 z_i=z_j\ \text{on every edge of }G.
\]
\end{theorem}

\begin{proof}
The \((i,j)\) entry of the matrix displacement is
\((z_i\overline z_j-1)B^{(a)}_{ij}\).  Squared Hilbert--Schmidt norm is the
sum of squared entry moduli, and summing over \(a\) gives the identity.
Every term is nonnegative.  A positive coefficient forces
\(z_i\overline z_j=1\), which is equivalent to \(z_i=z_j\) for unit phases.
The converse is immediate, including diagonal and zero-weight entries.
\end{proof}

\begin{proposition}[Observable phase rank]
\label{prop:rank}
The linear space of edge differences of \(G\) has real dimension
\[
 d-c(G).
\]
\end{proposition}

\begin{proof}
Choose one root in every connected component and a spanning tree within each
nontrivial component.  Differences from the component roots span every edge
difference, giving at most \(d-c(G)\) coordinates.  Conversely the
root-relative coordinates along the spanning forest are independent, giving
exactly \(d-c(G)\).
\end{proof}

\begin{theorem}[Graph-rank recurrence deadline]
\label{thm:deadline}
Let \(z_i(n)=\ee^{inE_i}\).  For every \(Q\in\N\), there is a positive
integer \(n\) such that
\[
 n\leq(Q+1)^{d-c(G)}
\]
and
\[
 D^2(z(n))\leq
 4\left(\frac{2\pi}{Q+1}\right)^2\sum_{i,j}W_{ij}.
\]
\end{theorem}

\begin{proof}
Choose one root in each component.  Simultaneous pigeonhole approximation
of the \(d-c(G)\) non-root relative energies gives
\[
 1\leq n\leq(Q+1)^{d-c(G)}
\]
and an integer \(m_i\) for every vertex such that
\[
 |n(E_i-E_{\operatorname{root}(i)})-2\pi m_i|
 \leq\frac{2\pi}{Q+1}.
\]
The error of an edge is the difference of two errors in the same component,
so it is at most \(2(2\pi/(Q+1))\).  Since
\[
 |\ee^{ix}-1|\leq|x-2\pi m|
\]
for every integer \(m\), each positively weighted edge term is at most the
square of that value.  Sum against \(W_{ij}\).
\end{proof}

\begin{corollary}[Operational probe and calibration bound]
\label{cor:operational}
Let
\(\Delta B^{(a)}=\operatorname{diag}(z)B^{(a)}
\operatorname{diag}(z)^\dagger-B^{(a)}\), assume \(\lambda_a>0\), and let
\(\sigma=\sigma^\dagger\) with \(\|\sigma\|_{\HS}\leq1\).  Then
\[
 |\operatorname{Tr}(\sigma\Delta B^{(a)})|
 \leq\sqrt{D^2(z)/\lambda_a}.
\]
If the final and initial calibrated matrices have Hilbert--Schmidt errors
bounded by \(\nu_1,\nu_0\), the measured expectation displacement is bounded
by
\[
 \sqrt{D^2(z)/\lambda_a}+\nu_1+\nu_0.
\]
\end{corollary}

\begin{proof}
The self-adjoint identity converts the trace to the Hilbert--Schmidt pairing.
Cauchy--Schwarz bounds its square by
\(\|\sigma\|_{\HS}^2\|\Delta B^{(a)}\|_{\HS}^2\).  The \(a\)th nonnegative
term of \(D^2\) is
\(\lambda_a\|\Delta B^{(a)}\|_{\HS}^2\).  Divide and take square roots.
For two calibration errors, apply the same normalized pairing bound to each
and use the triangle inequality.
\end{proof}

\section{Arbitrary pure-point spectra}

\begin{theorem}[Summable countable phase recurrence]
\label{thm:countable}
Let \(a_p\geq0\) with \(\sum_pa_p<\infty\), and let
\(\omega_p\in\R\) be arbitrary.  For
\[
 \mathcal E(n)=\sum_{p=0}^{\infty}a_p
   |\ee^{in\omega_p}-1|^2,
\]
every \(M,Q\in\N\) admit \(1\leq n\leq(Q+1)^M\) such that
\[
 \mathcal E(n)\leq
 \left(\frac{2\pi}{Q+1}\right)^2\sum_{p<M}a_p
 +4\sum_{p\geq M}a_p.
\]
Consequently, for every \(\varepsilon>0\), some \(n\geq1\) satisfies
\(\mathcal E(n)<\varepsilon\).
\end{theorem}

\begin{proof}
Simultaneous approximation of the first \(M\) frequencies produces
\(1\leq n\leq(Q+1)^M\) and head chords at most \(2\pi/(Q+1)\).
Every unit-circle chord is at most \(2\), which bounds the remaining squared
tail chords by \(4\).  This proves the quantitative inequality.  Given
\(\varepsilon>0\), choose \(M\) so the summable tail is below
\(\varepsilon/8\), then choose \(Q\) so the finite-head term is below
\(\varepsilon/2\).
\end{proof}

Let \(J\) and \(\iota\) be arbitrary index types, let \(B^{(a)}_{ij}\) lie
in a complex normed space, and set
\[
 \Phi_{ij}(t)=\ee^{it(E_i-E_j)}.
\]
Termwise,
\[
 \|\Phi_{ij}(t)B^{(a)}_{ij}-B^{(a)}_{ij}\|^2
 =\|B^{(a)}_{ij}\|^2|\Phi_{ij}(t)-1|^2.
\]
Therefore the observer displacement series is summable if and only if the
phase-energy series is summable, and their sums are equal.  Under
\(\lambda_a\geq0\), a summable total is zero if and only if every term with
\(\lambda_a>0\) and \(B^{(a)}_{ij}\ne0\) has \(\Phi_{ij}(t)=1\).  This proves
the arbitrary pure-point identity and zero law without a dimension
assumption.

\section{Explicit equal-information qutrit separation}

\smallskip
\noindent
Let
\[
 s=\sqrt2,\qquad \delta=s-1,\qquad
 p=(1-\delta,\delta,0),\quad
 q=(1-2\delta,\delta,\delta).
\]
Because \(1<s<3/2\), both packets are nonnegative and sum to one.  With
\(E=(0,1,s)\),
\begin{align*}
 \operatorname{Var}_{p}(E)
 &=\delta-\delta^2,\\
 \operatorname{Var}_{q}(E)
 &=3\delta-(1+s)^2\delta^2.
\end{align*}
Substitution of \(\delta=s-1\) and \(s^2=2\) makes these equal.

For a unit vector \(\psi\), the state
\[
 \rho_\psi=\frac16I+\frac12|\psi\rangle\langle\psi|
\]
has eigenvalue \(2/3\) along \(\psi\) and eigenvalue \(1/6\) on its
orthogonal complement.  Hence the two realized states have identical full
density spectra.  Direct differentiation of
\(\rho(t)=\ee^{-itH}\rho\ee^{itH}\) gives
\[
 \|\dot\rho(0)\|_{\HS}^2
 =2\left(\frac12\right)^2\operatorname{Var}(E).
\]
Writing \(a=2/3\), \(b=1/6\), and \(c=1/2\) for the spike eigenvalue,
background eigenvalue, and their difference, respectively, the
symmetric-logarithmic-derivative quantum Fisher information is
\[
 F_Q=\frac{4c^2}{a+b}\operatorname{Var}(E)
 =\frac65\operatorname{Var}(E).
\]
It is therefore also equal for the two states.

At \(t=2\pi n\), packet \(p\) sees only phases \(1\) and
\(\ee^{2\pi in}=1\), and therefore returns exactly.  Packet \(q\) also sees
\(\ee^{2\pi in\sqrt2}\).  Its recurrence amplitude is
\[
 R_q(n)=\frac12\sqrt{\delta(1-\delta)}
 \left|\ee^{2\pi in\sqrt2}-1\right|.
\]
For the nearest integer \(m\) to \(n\sqrt2\),
\[
 1\leq|2n^2-m^2|
 =|n\sqrt2-m|\,|n\sqrt2+m|
 \leq4n\,|n\sqrt2-m|,
\]
which, together with the circle chord bound, gives
\[
 R_q(n)\geq\frac{\sqrt{\delta(1-\delta)}}{2n}.
\]
Conversely, one-dimensional Dirichlet approximation supplies, for every
\(Q\geq1\), some \(1\leq n\leq Q\) whose circle distance is at most
\(2\pi/Q\), yielding
\[
 R_q(n)\leq\frac{\pi\sqrt{\delta(1-\delta)}}{Q}.
\]
Thus the second clock has a sharp inverse-resolution stroboscopic law even
though all listed local information agrees with the periodic clock.

\end{document}
